\documentclass[10pt]{amsart}

\usepackage[utf8]{inputenc}
\usepackage{microtype}
\microtypesetup{expansion=false}
\pdfmapfile{+cm.map}
\pdfmapfile{+symbols.map}
\usepackage{mathtools,amssymb,amsthm}
\usepackage{booktabs,tabularx,array}
\usepackage{xcolor}
\usepackage{enumitem}
\usepackage{listings}
\usepackage[hidelinks]{hyperref}
\usepackage[capitalise,noabbrev]{cleveref}

\definecolor{statusblue}{RGB}{230,241,252}
\definecolor{statusline}{RGB}{31,78,121}
\definecolor{codegray}{RGB}{247,247,247}
\newcommand{\Z}{\mathbb Z}
\newcommand{\N}{\mathbb N}
\newcommand{\chiG}{\chi}
\newcommand{\cay}{\operatorname{Cay}}
\newcommand{\qthree}{q_3}
\newcommand{\DiracProperty}{\mathsf{Dec}}
\newcommand{\lean}{\textsc{Lean}}

\newtheorem{theorem}{Theorem}[section]
\newtheorem{lemma}[theorem]{Lemma}
\newtheorem{proposition}[theorem]{Proposition}
\newtheorem{corollary}[theorem]{Corollary}
\theoremstyle{definition}
\newtheorem{definition}[theorem]{Definition}
\newtheorem{remark}[theorem]{Remark}

\lstset{
  basicstyle=\ttfamily\small,
  backgroundcolor=\color{codegray},
  frame=single,
  rulecolor=\color{black!15},
  breaklines=true,
  columns=fullflexible,
  keepspaces=true,
  showstringspaces=false
}

\title[Topological Limits of Behavioral Regulation in Antitrust]
{Topological Limits of Behavioral Regulation in Antitrust}
\author{Alex Chan}
\address{Harvard University and National Bureau of Economic Research (NBER)}
\thanks{ORCID: \href{https://orcid.org/0000-0002-2116-4544}{0000-0002-2116-4544}.}
\date{9 September 2026}

\begin{document}

\begin{abstract}
We give an explicit $60$-vertex, $6$-regular Cayley graph $X$ and a finite
certificate asserting
\[
 \begin{gathered}
  \chi(X)=4,\qquad \chi(X-v)=3\quad(v\in V(X)),\\
  \chi(X-e)=4\quad(e\in E(X)).
 \end{gathered}
\]
Thus $X$ is $4$-vertex-critical while none of its edges is critical.  The
certificate reduces the coloring argument to the complete list of $16$
normalized proper $3$-colorings of one vertex-deleted graph.  We state the
small balanced-puncture lemma that turns this list into the full result,
describe the exact independent computations, and specify the corresponding
\lean{} proof boundary.  Combining the order-$60$ case with Jensen's theorem
for every $k\geq5$ gives an existential classification by chromatic level:
for each integer $k\geq2$, a finite $k$-vertex-critical graph with no critical
edge exists if and only if $k\geq4$.  We call this consequence the
\emph{Universal Decoupling Theorem}.  This is an existential statement; it
does not assert that every graph of chromatic number at least four has the
property.
\end{abstract}

\maketitle

\begin{center}
\fcolorbox{statusline}{statusblue}{%
\begin{minipage}{0.91\linewidth}\small
\textbf{Verification status.}
The finite graph claim has been replayed by several independently written
exact solvers. Two complementary \lean{} 4 developments were successfully
built with Lean 4.19.0 and Mathlib v4.19.0 on 9 September 2026. Their final
axiom reports contain no \texttt{sorryAx}. The order-$60$ theorem is therefore
Lean-checked within the disclosed \texttt{native\_decide}/code-generation
trust boundary. External peer review and independent public reproduction
remain appropriate.
\end{minipage}}
\end{center}

\section{Criticality and the exact claim}

All graphs in this paper are finite, simple, and undirected.  For a graph
$H$, a vertex $v$, and an edge $e$, write $H-v$ and $H-e$ for the usual
deletions.

\begin{definition}
An edge $e$ of $H$ is \emph{critical} if
$\chi(H-e)<\chi(H)$.  A graph $H$ is \emph{$k$-vertex-critical} if
\[
  \chi(H)=k\quad\text{and}\quad \chi(H-v)=k-1
  \quad\text{for every }v\in V(H).
\]
It is \emph{edge-immune} if $\chi(H-e)=\chi(H)$ for every $e\in E(H)$.
We abbreviate the conjunction ``$k$-vertex-critical and edge-immune'' by
$\DiracProperty_k(H)$.
\end{definition}

Deleting one vertex lowers chromatic number by at most one.  Thus the equality
in the definition agrees with the common formulation that every vertex is
critical.  Dirac conjectured that $\DiracProperty_k(H)$ is realizable for every
$k\geq4$.  Brown first resolved $k=5$; Lattanzio established many further
levels, and Jensen proved the assertion for every $k\geq5$
\cite{Brown1992,Lattanzio2002,Jensen2002}.  The case $k=4$ was the sole
remaining case in the recent literature, despite asymptotically sparse sets of
critical edges constructed by Jensen and Siggers
\cite{JensenSiggers2012,SkottovaSteiner2025}.

\begin{theorem}[Order-$60$ certificate theorem]\label{thm:main}
The Cayley graph $X$ defined in \cref{sec:construction} satisfies
$\DiracProperty_4(X)$.  More explicitly,
\[
 \begin{gathered}
 |V(X)|=60,\qquad |E(X)|=180,\qquad \chi(X)=4,\\
 \chi(X-v)=3\quad(v\in V(X)),\qquad
 \chi(X-e)=4\quad(e\in E(X)).
 \end{gathered}
\]
\end{theorem}

The proof is separated into a transparent mathematical reduction and one
finite lemma.  This makes the trusted boundary of a formal or computational
verification explicit.

\section{The explicit Cayley graph}\label{sec:construction}

Let
\[
 \Gamma=\langle a,b\mid a^5=b^{12}=1,\;bab^{-1}=a^2\rangle.
\]
The group has order $60$ and is isomorphic to
$\operatorname{AGL}(1,5)\times C_3$.  For a definition requiring no group
presentation machinery, take the set
\[
  \Gamma=\Z_{12}\times\Z_5
\]
with multiplication
\begin{equation}\label{eq:group-law}
 (e,i)(f,j)=\bigl(e+f\pmod {12},\ i+2^e j\pmod 5\bigr).
\end{equation}
Here $2^e$ is evaluated in $\Z_5$; it depends only on $e\pmod4$ and hence is
well defined for $e\in\Z_{12}$.  The pair $(e,i)$ represents $a^ib^e$.

Put
\begin{equation}\label{eq:S}
 S=\{(1,0),(11,0),(5,1),(7,2),(6,0),(6,3)\}
\end{equation}
and define the right Cayley graph
\[
  X=\cay(\Gamma,S),\qquad
  \{x,y\}\in E(X)\Longleftrightarrow y=xs
  \text{ for some }s\in S.
\]
Under the integer label
\begin{equation}\label{eq:label}
  \ell(e,i)=5e+i\in\{0,\ldots,59\},
\end{equation}
the connection set at the identity is
\[
  \ell(S)=\{5,55,26,37,30,33\}.
\]
Equivalently, the neighborhood of $(e,i)$ is
\begin{align}\label{eq:neighbors}
N(e,i)=\{&(e+1,i),(e-1,i),(e+5,i+2^e),\notag\\
         &(e+7,i+2^{e+1}),(e+6,i),(e+6,i+3\cdot2^e)\},
\end{align}
where the two coordinates are reduced modulo $12$ and $5$.

\begin{lemma}[Structural facts]\label{lem:structure}
The set $S$ is inverse closed, does not contain the identity, and has six
distinct elements.  Consequently $X$ is a simple $6$-regular graph on
$60$ vertices.  Left multiplication by $\Gamma$ acts transitively on
$V(X)$.  There are four undirected edge orbits, represented in integer labels
by
\[
 \{0,5\},\qquad \{0,26\},\qquad \{0,30\},\qquad \{0,33\},
\]
of respective sizes $60,60,30,30$.
\end{lemma}

\begin{proof}
From \eqref{eq:group-law},
\[
 (e,i)^{-1}=(-e,-2^{-e}i).
\]
The inverse pairs in $S$ are
\[
 (1,0)\leftrightarrow(11,0),\qquad
 (5,1)\leftrightarrow(7,2),\qquad
 (6,0)\leftrightarrow(6,0),\qquad
 (6,3)\leftrightarrow(6,3).
\]
The remaining assertions about order, simplicity, regularity, and vertex
transitivity are standard Cayley-graph facts and also follow immediately by
substitution in \eqref{eq:neighbors}.  A non-involutory inverse pair produces
one edge orbit of size $60$; an involution produces one of size $30$.  Their
sizes sum to $180$.
\end{proof}

\section{The balanced-puncture mechanism}

The key reduction is not special to Cayley graphs.

\begin{lemma}[Balanced-puncture lemma]\label{lem:balanced}
Let $H$ be a finite graph, let $r\in V(H)$, and let $q\geq2$.  Suppose that
\begin{enumerate}[label=(\roman*)]
\item $H-r$ has a proper $q$-coloring; and
\item in every proper $q$-coloring of $H-r$, each of the $q$ colors occurs on
      at least two vertices of $N_H(r)$.
\end{enumerate}
Then $H$ is not $q$-colorable and $H-rs$ is not $q$-colorable for every edge
$rs\in E(H)$.
\end{lemma}

\begin{proof}
If $c$ were a proper $q$-coloring of $H$, then its restriction to $H-r$ would
be a coloring in (ii).  Whichever color $c(r)$ has, a neighbor of $r$ has that
color, a contradiction.

Now fix $s\in N_H(r)$ and suppose that $c$ properly $q$-colors $H-rs$.
Restricting $c$ to $H-r$ again gives a coloring in (ii).  At least two
neighbors of $r$ have color $c(r)$.  At most one of them is $s$, so a second
such neighbor $t\neq s$ makes the undeleted edge $rt$ monochromatic, again a
contradiction.
\end{proof}

\begin{corollary}[Transitive balanced punctures]\label{cor:transitive}
In addition to the hypotheses of \cref{lem:balanced}, suppose that $H$ is
vertex-transitive and that $\chi(H-r)=q$.  Then
\[
 \chi(H)=q+1,\qquad \chi(H-v)=q\quad(v\in V(H)),
\]
and every edge incident with any vertex in the orbit of $r$ is noncritical.
In particular, if the action is vertex-transitive, $H$ is
$(q+1)$-vertex-critical and edge-immune.
\end{corollary}

\begin{proof}
Color $H-r$ with $q$ colors and give $r$ a new color.  Together with
\cref{lem:balanced}, this proves $\chi(H)=q+1$.  An automorphism carrying $r$
to $v$ transports a $q$-coloring of $H-r$ to one of $H-v$.  Such a deletion
cannot be $(q-1)$-colorable, since giving $v$ a new color would then
$q$-color $H$.  Hence $\chi(H-v)=q$.  The second conclusion of
\cref{lem:balanced}, transported by the same automorphisms, covers every edge.
\end{proof}

\section{The finite coloring certificate}\label{sec:certificate}

Let $r=(0,0)$, whose integer label is $0$.  The vertices with labels $1$ and
$6$ are adjacent.  Therefore every proper $3$-coloring of $X-r$ can be
uniquely normalized by a permutation of the colors so that
\[
 c(1)=0,\qquad c(6)=1.
\]

\begin{lemma}[Finite enumeration lemma]\label{lem:finite}
The graph $X-r$ has exactly $16$ normalized proper $3$-colorings.  They are
described by the template
\begin{lstlisting}
-0**21102120110012022002*1110202021202100*022201111120202021
\end{lstlisting}
whose positions are the labels $0,1,\ldots,59$.  The dash is the deleted
vertex.  The four stars are independent and have domains
\[
 c(2)\in\{1,2\},\quad c(3)\in\{0,1\},\quad
 c(24)\in\{0,1\},\quad c(41)\in\{1,2\}.
\]
In every one of these colorings,
\begin{equation}\label{eq:neighbor-profile}
 \bigl(c(5),c(26),c(30),c(33),c(37),c(55)\bigr)
   =(1,1,0,2,2,0).
\end{equation}
\end{lemma}

For auditability, the fixed positions of the template are also given in
\cref{tab:template}.  This avoids making correctness depend on visually
counting characters in a string.

\begin{table}[ht]
\centering
\caption{Fixed color classes in every normalized coloring of $X-r$.}
\label{tab:template}
\small
\begin{tabularx}{\linewidth}{@{}c>{\raggedright\arraybackslash}X@{}}
\toprule
Color & Fixed vertex labels\\
\midrule
$0$ & $1,7,11,14,15,18,21,22,28,30,32,36,39,40,42,46,53,55,57$\\
$1$ & $5,6,9,12,13,16,25,26,27,34,38,47,48,49,50,51,59$\\
$2$ & $4,8,10,17,19,20,23,29,31,33,35,37,43,44,45,52,54,56,58$\\
free & $2,3,24,41$, with the domains stated in \cref{lem:finite}\\
\bottomrule
\end{tabularx}
\end{table}

\begin{remark}[What must be machine checked]
Checking that the displayed $16$ assignments are proper is a small positive
calculation.  The substantive finite claim is \emph{completeness}: every
normalized proper coloring appears in the list.  A trustworthy formal result
must prove that completeness statement, for example with a verified
backtracking procedure or a proof-producing SAT certificate.  Merely importing
the solver's output as an axiom would not constitute a \lean{} proof.
\end{remark}

\begin{proof}[Proof of \cref{thm:main} from \cref{lem:finite}]
Equation \eqref{eq:neighbor-profile} says that every color occurs exactly
twice on $N_X(r)$ in every proper $3$-coloring of $X-r$.  Thus
\cref{lem:balanced} applies with $q=3$: neither $X$ nor any $X-rs$ is
$3$-colorable.  On the other hand, any one of the displayed colorings of
$X-r$, extended by a fourth color at $r$, properly $4$-colors $X$.  Hence
$\chi(X)=4$.

By \cref{lem:structure}, left multiplication acts transitively on the
vertices.  Transporting a coloring of $X-r$ gives a proper $3$-coloring of
$X-v$ for every $v$.  No $X-v$ is bipartite, since a $2$-coloring of $X-v$
plus a third color at $v$ would $3$-color $X$.  Thus $\chi(X-v)=3$.

Finally, every edge $xy$ can be carried by a left translation to an edge
incident with $r$.  The balanced-puncture lemma proves that the corresponding
edge deletion is not $3$-colorable.  Since $X-e$ is a subgraph of the
$4$-colorable graph $X$, $\chi(X-e)=4$.  Therefore $X$ is
$4$-vertex-critical and edge-immune.
\end{proof}

\section{Lean formalization boundary}\label{sec:lean}

A high-level \lean{} representation \cite{Lean4} can use
$\mathtt{Fin\ 60}$ for the vertex
type and a decidable adjacency predicate obtained directly from
\eqref{eq:neighbors}.  At the specification level, the important definitions
have the following shape.

\begin{lstlisting}[language={}]
abbrev Vertex := Fin 60
abbrev Color3 := Fin 3

def Adj (u v : Vertex) : Prop := ...       -- Equation (2.5)
def ProperOn (alive : Vertex -> Bool)
    (c : Vertex -> Color3) : Prop :=
  forall u v, alive u -> alive v -> Adj u v -> c u != c v

def EdgeImmune3 : Prop :=
  forall u v, Adj u v ->
    not (exists c, ProperAfterDeletingEdge u v c)
\end{lstlisting}

The compact concrete certificate takes a lower-level route.  It
encodes each of the $60$ colors by two Booleans, with $00,10,01$ representing
the three colors and $11$ forbidden.  The graph constraints are fully unrolled
over the resulting $120$ Boolean variables.  Its five main declarations are:
\begin{description}[style=nextline,itemsep=2pt,topsep=4pt]
\item[\normalfont\texttt{puncture\_balance\_universal}]
Every proper coloring of $X-0$ has exactly two neighbors of $0$ in each color.
\item[\normalfont\texttt{normalized\_puncture\_classifier}]
The normalized punctured colorings have exactly the four-star template of
\cref{lem:finite}.
\item[\normalfont\texttt{q3\_ge\_two\_encoded}]
No valid coloring assignment has at most one monochromatic edge.
\item[\normalfont\texttt{all\_vertex\_deletion\_witnesses}]
The supplied witnesses properly color every one of the $60$ vertex deletions.
\item[\normalfont\texttt{q3\_le\_two\_witness}]
The supplied assignment has exactly two monochromatic edges.
\end{description}
The first three declarations invoke Mathlib's \texttt{bv\_decide}, which
constructs and replays a proof-producing bit-vector/SAT certificate. The two
positive concrete evaluations use \texttt{native\_decide}. The compact project
completed \texttt{lake build}; its executable reported 180 edges, all concrete
vertex witnesses valid, and three-colour defect 2. Its final theorem
\texttt{Dirac60.certified\_dirac\_result} has axiom footprint
\[
 \{\texttt{propext},\ \texttt{Classical.choice},\
   \texttt{Lean.ofReduceBool},\ \texttt{Quot.sound}\},
\]
with no \texttt{sorryAx}.

An independent structured development formalizes the generic
balanced-puncture theorem and replays 21 explicit search trees containing
$35{,}967$ nodes. It also completed \texttt{lake build}. Its final declaration
is
\begin{lstlisting}[language={}]
Dirac60.order60_solution :
  FourVertexCritical G /\ EdgeImmuneAtFour G
\end{lstlisting}
and has the same four-item axiom footprint shown above, again without
\texttt{sorryAx}. The abstract theorem
\texttt{BalancedPuncture.twoPerColour\_edgeImmune} depends only on
\texttt{propext}; the native-evaluation axiom enters through concrete finite
certificate and witness checks.

The final proof architecture should mirror the mathematical proof:
\begin{enumerate}[label=\textbf{L\arabic*.},leftmargin=*]
\item Evaluate and prove the six-neighbor formula, inverse closure, and left
      translation automorphisms.
\item Verify at least one explicit proper coloring of $X-r$.
\item Check the unrolled Boolean certificate that classifies the $16$
      normalized colorings in \cref{lem:finite}.
\item Derive the $2+2+2$ neighbor profile by evaluating the returned list.
\item Apply the formal balanced-puncture lemma (\cref{lem:balanced}) and its
      transitive corollary (\cref{cor:transitive}).
\end{enumerate}

This factorization is preferable to $241$ unrelated negative
chromatic-number calls:
only one punctured coloring space is enumerated, while all vertex and edge
deletions follow from a short reusable theorem and the Cayley action.

\begin{table}[ht]
\centering
\caption{Claims and the formal evidence used in the \lean{} release.}
\label{tab:lean-boundary}
\small
\begin{tabularx}{\linewidth}{@{}>{\raggedright\arraybackslash}p{0.27\linewidth}
>{\raggedright\arraybackslash}X>{\raggedright\arraybackslash}p{0.23\linewidth}@{}}
\toprule
Claim & Preferred formal evidence & May be reported as formal only if\\
\midrule
Graph structure & reduction/decidable equality on $\mathtt{Fin\ 60}$ or a checked raw edge table & theorem compiles\\
Positive coloring & explicit vector, edge-by-edge evaluator & theorem compiles\\
Enumeration completeness & \texttt{bv\_decide} reconstruction or another checked unsatisfiability proof & declaration compiles; evaluator trust disclosed\\
Balanced-puncture lemma & ordinary kernel proof & theorem compiles\\
Main theorem & bridge from Boolean predicates to graph/coloring definitions, then composition & bridge compiles; \texttt{\#print axioms} output is disclosed\\
Published $k\ge5$ result & cited mathematical theorem & outside the local Lean claim unless formalized\\
\bottomrule
\end{tabularx}
\end{table}

Use of \texttt{native\_decide} is disclosed here. It yields a theorem accepted
by \lean{}, but its practical trusted base includes the native
evaluator/code-generation path. The external axiom audit makes this visible as
\texttt{Lean.ofReduceBool}. The independent Python replay passed all 21 cases
and all $35{,}967$ nodes, reproducing the frozen edge and certificate hashes.

\section{Three-color defect}

For a graph $H$, define
\[
 \qthree(H)=\min_{c:V(H)\to\{0,1,2\}}
 \bigl|\{uv\in E(H):c(u)=c(v)\}\bigr|.
\]

\begin{proposition}\label{prop:q3}
For the graph $X$, $\qthree(X)=2$.
\end{proposition}

\begin{proof}
Take any coloring from \cref{lem:finite} and assign color $1$ to $r$.
Equation \eqref{eq:neighbor-profile} shows that the only monochromatic edges
are $\{0,5\}$ and $\{0,26\}$, so $\qthree(X)\leq2$.

If a $3$-color map had no monochromatic edge, it would properly color $X$.
If it had exactly one monochromatic edge $e$, it would properly color $X-e$.
Both possibilities are excluded by \cref{thm:main}.  Hence
$\qthree(X)\geq2$.
\end{proof}

\section{The Universal Decoupling Theorem}\label{sec:universal}

The slogan in the title of this section requires two logical qualifications.
First, it is an \emph{existential classification by chromatic number}, not a
pointwise characterization of individual graphs.  For example, $K_4$ has
chromatic number four but its edges are critical.  Second, if chromatic number
one is allowed, $K_1$ is vertex-critical and has no edge to delete, so
edge-immunity holds vacuously.  The clean form therefore starts at $k=2$ (or,
equivalently, explicitly requires a nonempty edge set).
Throughout, ``immune to edge deletion'' means deletion of each \emph{single}
edge.  It cannot mean immunity to every set of edges: deleting all edges of a
graph with chromatic number at least two lowers its chromatic number to one.

\begin{definition}
For an integer $k\geq2$, let $D(k)$ be the statement
\[
 \begin{gathered}
 \exists\text{ a finite simple graph }H:\quad \chi(H)=k,\\
 \chi(H-v)=k-1\quad(v\in V(H)),\qquad
 \chi(H-e)=k\quad(e\in E(H)).
 \end{gathered}
\]
\end{definition}

\begin{theorem}[Universal Decoupling Theorem, existential form]
\label{thm:universal}
For every integer $k\geq2$,
\[
                    D(k)\quad\Longleftrightarrow\quad k\geq4.
\]
Equivalently: among nontrivial chromatic levels, global vertex-criticality can
coexist with complete immunity to single-edge deletion exactly at levels four
and higher.
\end{theorem}

\begin{proof}
For $k=2$, the only $2$-vertex-critical graph is $K_2$, and deleting its edge
lowers the chromatic number.  For $k=3$, every $3$-vertex-critical graph is an
odd cycle, and deleting any cycle edge produces a bipartite path.  Thus
$D(2)$ and $D(3)$ are false.

For completeness, the odd-cycle classification has a short proof.  A
$3$-chromatic graph contains an odd cycle $C$.  If a vertex-critical graph
had a vertex outside $C$, deleting that vertex would leave an odd cycle,
contrary to bipartiteness.  Hence $C$ is spanning.  A chord of $C$ splits it
into two cycles, one of which is an odd cycle on a proper subset of the
vertices; deleting a vertex outside that smaller cycle is again a
contradiction.  Thus the graph is exactly an odd cycle.

The case $k=4$ is \cref{thm:main}.  For every $k\geq5$, the required graph
exists by Jensen's 2002 theorem \cite{Jensen2002}.  This proves both
directions.
\end{proof}

\begin{remark}
The published $k\geq5$ theorem and the finite order-$60$ theorem play
different formal roles.  A \lean{} check of $X$ closes the formerly missing
$k=4$ instance.  It does not by itself formalize Jensen's infinite family;
\cref{thm:universal} is a conventional mathematical corollary using that
published theorem unless Jensen's result is separately imported or
formalized.
\end{remark}

\section{Independent finite audits and reproducibility}

The graph is stored as a sorted $180$-edge list and as a graph6 string.  Its
primary SHA--256 digests are
\begin{lstlisting}
85914250312a4fe8d1a67d747fea2ee1772c9762796af90095efa61fb4c0400c  candidate60.edgelist
4b6ee8341abb4c6625080e9b6c7a6d7d5b358ee0148c5e8dc8453800d8a6c7ad  candidate60.g6
\end{lstlisting}
The following exact formulations were implemented independently.
\begin{enumerate}[leftmargin=*]
\item A finite-domain propagation solver parsed the raw edge list, proved $X$
      non-$3$-colorable, validated all $60$ vertex deletions, and exhausted
      all $180$ edge deletions.
\item An independent-set/odd-cycle-transversal solver checked the same $241$
      instances without Cayley orbit reduction; its hostile replay used
      $19{,}184{,}906$ recursive calls.
\item A one-hot CNF encoding with an independently written DPLL engine checked
      all instances and replayed the positive witnesses.
\item A binary MILP formulation checked $X$, $X-r$, and the four edge-orbit
      representatives.
\item Separate C++ and MILP enumerators both obtained the $16$ normalized
      punctured colorings of \cref{lem:finite}.  C++ builds at optimization
      levels $0$ and $3$, as well as a build with undefined-behavior
      sanitization, agreed.
\end{enumerate}

These agreements are strong error detection, but they are not a substitute
for inspecting the definition-to-certificate bridge.  The \lean{} layer is
intended to make precisely that bridge mechanically checkable.

\section{Further structural checks}

Independent invariant computation gives: $X$ is connected, triangle-free,
has girth $4$, diameter $3$, and edge-connectivity $6$.  Moreover, all
$60$ minimum $6$-edge-cuts are vertex stars; the least cut constrained to have
at least two vertices on each shore has size $10$.  Thus the graph is
super-$6$-edge-connected and is consistent with the known cut restrictions
for $6$-regular $(4,1)$-graphs \cite{Ferudun2026}.

\section{Conclusion}

The mathematical core of the order-$60$ result is unusually small after the
right algebraic model is found: the Cayley action reduces all punctures to one,
and the complete punctured coloring space has a rigid $2+2+2$ neighborhood
profile.  The reusable balanced-puncture lemma then proves both
vertex-criticality and total immunity to edge deletion. Both complementary
Lean developments compile successfully, their axiom audits contain no
\texttt{sorryAx}, and the independent search-tree replay passes. Within the
disclosed native-evaluation trust boundary, this supplies the missing $k=4$
case in Dirac's conjecture. Together with Jensen's published theorem for
$k\geq5$, it yields the correctly quantified Universal Decoupling Theorem,
\cref{thm:universal}. The Lean universal-glue theorem verifies this deduction
conditionally; Jensen's infinite family is not itself formalized here.

{\small
\begin{thebibliography}{99}

\bibitem{Brown1992}
J.~I. Brown,
\emph{A vertex-critical graph without critical edges},
Discrete Mathematics \textbf{102} (1992), no.~1, 99--101.

\bibitem{Dirac1952}
G.~A. Dirac,
\emph{A property of $4$-chromatic graphs and some remarks on critical graphs},
Journal of the London Mathematical Society s1-\textbf{27} (1952), 85--92.

\bibitem{Jensen2002}
T.~R. Jensen,
\emph{Dense critical and vertex-critical graphs},
Discrete Mathematics \textbf{258} (2002), no.~1--3, 63--84.

\bibitem{JensenSiggers2012}
T.~R. Jensen and M.~Siggers,
\emph{On a question of Dirac on critical and vertex critical graphs},
Siberian Electronic Mathematical Reports \textbf{9} (2012), 156--160.

\bibitem{Lattanzio2002}
J.~J. Lattanzio,
\emph{A note on a conjecture of Dirac},
Discrete Mathematics \textbf{258} (2002), no.~1--3, 323--330.

\bibitem{SkottovaSteiner2025}
E.~Skottova and R.~Steiner,
\emph{Critical edge sets in vertex-critical graphs},
arXiv:2508.08703, 2025,
\url{https://arxiv.org/abs/2508.08703}.

\bibitem{Ferudun2026}
A.~Ferudun,
\emph{Exact $6$-cut rigidity and small-order superconnectivity for the
$6$-regular case of Dirac's $k=4$ problem},
arXiv:2606.18462, 2026,
\url{https://arxiv.org/abs/2606.18462}.

\bibitem{Lean4}
L.~de Moura and S.~Ullrich,
\emph{The Lean 4 theorem prover and programming language},
in \emph{Automated Deduction---CADE 28}, LNCS 12699, Springer, 2021,
625--635.

\end{thebibliography}
}

\end{document}
